\chapter{结论与展望}
\section{主要结论}
本研究围绕示例研究课题开展了系统的研究工作，主要结论如下：

（1）研究方法验证：本研究提出的数值分析方法经过试验验证，能够准确模拟研究对象的关键力学性能，误差控制在5\%以内，表明研究方法具有较高的可靠性。

（2）参数影响规律：通过参数分析，揭示了各关键参数对研究结果的影响规律。屈服强度对结果影响最为显著，相对变化达到20\%；弹性模量影响中等，相对变化约15\%；几何参数影响相对较小。

（3）机理揭示：通过数值分析和试验验证，揭示了研究对象在不同工况下的力学机理，包括荷载传递机理、变形机理和破坏机理。

（4）工程应用价值：基于研究结果提出了设计建议和施工建议，为实际工程应用提供了参考依据。

\section{创新点}
本研究的创新点主要包括：
\begin{enumerate}
    \item 提出了一种新的数值分析方法，能够准确模拟研究对象的复杂力学行为。
    \item 系统揭示了关键参数对研究结果的影响规律，填补了相关研究空白。
    \item 提出了面向工程应用的设计建议，具有实际应用价值。
\end{enumerate}

\section{研究展望}
尽管本研究取得了一定成果，但仍有以下问题值得进一步研究：

（1）研究范围拓展：本研究仅针对特定工况进行研究，建议拓展研究范围，验证结论的普适性。

（2）方法改进：数值分析方法存在简化假设，建议进一步完善分析模型，提高结果的准确性。

（3）工程验证：建议开展更多的工程验证研究，检验研究成果在实际工程中的应用效果。

（4）新方法探索：建议探索新的研究方法，如机器学习方法等，提高研究效率。